% 00-about.tex: what this document is, and how to read its marks.
\section*{About this version}
\addcontentsline{toc}{section}{About this version}

Version 1 of this report (\code{polaris\_project\_report.pdf}, preserved unchanged beside this
file) described Polaris as a database-design project: twelve tables around one credential, an
entity-relationship model, a lifecycle state machine, and the argument that identity data must be
post-quantum from the first day. Version 2 (\code{polaris\_project\_report\_v2.pdf}, also preserved)
described the system as it stood at tag \code{v9.345}, thirty-six tables, a federation protocol, a
transparency layer and an institutional exchange fabric, and was revised at \code{v9.355} when the
holder-side work it had recorded as missing was built. Version 3 describes the system at version
\code{1.0.0-rc.7} (commit \code{4e74376}, 20 September 2026), re-measured whole rather than patched,
and it is written for a reader who has never opened the repository. It was first written at
\code{1.0.0-rc.1} four days earlier and re-measured in place; fourteen of the twenty-six counts in
Appendix~\ref{app:counts} did not move, and the ones that did are marked there.

\textbf{What changed, in one paragraph.} Between Version 2 and this one the repository stopped
counting its own progress and started counting witnesses. The four standalone packages a relying
party installs are on public registries. A wallet written by somebody else, run unmodified,
presented a credential to the verifier and was accepted, after refusing it once for a defect every
internal instrument had passed over. The OpenID Foundation's hosted conformance suite ran the
verifier profile across the open internet. The verification layer itself was attacked from inside:
the checks, the constraints, the triggers, the unique indexes, the stored procedures, the zero-knowledge
witnesses, the conformance contract and both reference verifiers were each mutated and required to
go red, and many did not, at first. The physical layer became an object somebody else can implement.
And two words on the front door were weakened, \emph{post-quantum} to \emph{algorithm agility under an
audited migration path} and \emph{compulsion-resistant} to \emph{duress-aware}, because the evidence
did not support the stronger ones. Appendix~\ref{app:from-v2} lists every claim of Version 2 that this
version withdraws or narrows.

\textbf{Three dates, stated plainly.} Version 2's body was measured at \code{v9.345}; its
holder-side sections at \code{v9.355}. Everything in this version, every count, every mark and every
figure, was measured at \code{1.0.0-rc.7} by the methods in Appendix~\ref{app:counts}. Where a number
is older than that, because the repository has not re-run the instrument that produced it, the paper
says so beside the number. Nothing that ran at \code{v9.355} stopped running.

The paper keeps what Version 1 got right, the token-centred geometry of the schema, used here as the
centre of every map, and the lifecycle machine, and what Version 2 got right, the concentric reading
of the system, one ring per section, each ring introduced by the failure of the ring inside it.

\subsection*{Five marks}

Every capability carries one of five marks, so that code that runs is never confused with an
intention, and so that a thing the author has tested is never confused with a thing somebody else
has.

\begin{description}
  \item[\sImpl] A path that runs at \code{1.0.0-rc.7}, exercised by a test or a drill in continuous
    integration, and pinned by an invariant check that fails the build if it stops being true.
  \item[\sExp] A path that runs but is limited in a way the repository itself names: not covered by
    CI, notional in scale, or dependent on a backend that CI simulates.
  \item[\sExt] A property the repository implements and tests internally but that no external
    party has audited, attacked, deployed or independently re-implemented.
  \item[\sWit] A path that a named party outside the repository has actually exercised, on a date,
    with a result that can be pointed at: a wallet, a conformance service, a published test corpus.
    The mark is deliberately rare, it is given only where the repository's scoreboard has a filled
    row, and it never means audited, certified or piloted. Three things carry it in this paper.
  \item[\sFut] A proposal: a roadmap row, a design note, or an idea that exists only in this paper.
    It is not code.
\end{description}

A row carrying two marks runs today and still requires external validation for the stronger claim,
or has been exercised from outside on one path and not on the others.

Numbers were measured on the tree at \code{1.0.0-rc.7} by the methods in
Appendix~\ref{app:counts}, not copied from earlier documents. Where the repository's own documents
disagree with its code, the code wins and the disagreement is recorded in the author's notes that
accompany this version.

\subsection*{The analogy, and its limit}

The paper explains Polaris through one sustained analogy: a walled town. The credential record is
the town hall's register; the constraints are the wall; the endpoints are gates; the transparency
logs are watchtowers; independent witnesses stand outside the wall watching the towers; other
authorities are other towns, joined by roads that run one way; above the whole town stand laws that
bind even the town itself; and, new in this version, travellers from elsewhere have now come to the
gate, and what they found there is the only report the town did not write about itself. The analogy
describes arrangement and sight lines, not political centralisation: Polaris is federated per
authority, with no central database, root or verification service, and every recurrence of the
analogy returns to the actual mechanism by its real name.

\subsection*{One thesis, and one rule}

The system is an attempt to build identity infrastructure in which important claims do not have to
be trusted merely because a central application says they are true. Authority, state, history,
privacy boundaries and the system's own behaviour are progressively exposed to independent
verification and constrained by layers that do not share each other's assumptions. Each layer exists
because the layer inside it could be entirely correct and still fail to prove what a society needs
to know. That progression is the spine of the paper.

The rule, adopted on 13 September 2026 and applied to everything since, is that the unit of progress
is a new external dependency survived: a named wallet, a named conformance profile, a named relying
party, a defect filed by somebody who is not the author. The count of invariants, versions and lines
is not progress and the repository's scoreboard refuses to record it. This paper follows the rule.
Where it reports internal evidence, it says so, and it keeps that evidence out of the sentences that
report the other kind.
